\chapter{Data: origin, meaning, and the synthetic-data method}
\label{ch:data}

\begin{quote}\itshape
This chapter in brief: before any method is judged, the reader deserves to know exactly
what the data \emph{is} --- what an airline actually measures, what each recorded number
physically represents, what ``synthetic data'' means as a scientific instrument, and the
complete provenance chain from certified test-bed measurements to the rows the algorithms
consume. Nothing downstream makes sense without this chapter; nothing in it requires more
than \cref{ch:background}.
\end{quote}

\keyidea{Real fleets offer no ground truth and almost no labels, so attribution --- did a method blame the right component? --- is unmeasurable on them. Synthetic-with-truth is not a compromise but the only instrument that can measure it; self-consistency (exact scoring) and physical fidelity (an audited question) are kept strictly separate.}


\section{What engine-health data exists in the real world}
\label{sec:real-data}

An airline's engine-monitoring data ecosystem has three tiers, and knowing them explains
every design choice in SynCFM56:

\begin{description}
  \item[ACARS snapshot reports (the tier this project mimics).] The Aircraft
    Communications Addressing and Reporting System transmits short, structured telegrams
    from aircraft to ground over VHF or satellite. For engine monitoring, the avionics
    capture a \emph{stabilized frame} --- a few seconds of averaged sensor values once
    per flight phase --- typically one report at takeoff and one in cruise. Each report
    carries on the order of 15--30 parameters per engine. This is the data tier on which
    classical EHM (and most airline practice) actually runs: \textbf{two rows per engine
    per flight}, not continuous signals. Its economics shaped it: ACARS transmission
    historically cost money per kilobyte, and a stabilized snapshot compresses a flight
    into exactly what trend monitoring needs.
  \item[Continuous recordings (QAR/CEOD).] The Quick Access Recorder and its modern
    full-flight equivalents record hundreds of parameters at 1--16\,Hz for the whole
    flight. Rich, but heavy: gigabytes per aircraft-day, often not downloaded every
    flight, historically processed only after events. Modern ``big data'' EHM programs
    increasingly exploit this tier; most published fleet experience still lives in tier
    one.
  \item[Shop and maintenance records.] Borescope findings, module performance at overhaul,
    parts scrapped, wash logs. This is where \emph{ground truth} about component condition
    is glimpsed --- but only at maintenance events, only qualitatively, and recorded in
    systems never designed for machine learning.
\end{description}

Two structural facts follow. First, \textbf{real fleets provide almost no usable ground
truth}: between shop visits (years apart), nobody knows the true efficiency of an
installed HPT to a tenth of a percent --- there is no oracle to score a diagnosis
against. Second, \textbf{failures are rare and unlabeled}: a dataset of ten thousand
real flights may contain zero acute gas-path faults, and when one occurs, its onset time
and root cause are reconstructed after the fact, imperfectly. These two facts --- no
truth, no labels --- are why this project generates its data, and they define exactly
what the generation must provide.

\section{What each measured channel physically represents}
\label{sec:channel-meaning}

The eight channels recorded in every SynCFM56 snapshot are the industry-standard gas-path
set. What each number \emph{is}, physically, and why monitoring cares:

\begin{description}
  \item[N1 {[}rpm{]} --- fan / low-pressure spool speed.] The big front rotor. On the
    CFM56-7B it is also the \emph{thrust-setting parameter}: the crew (or FADEC) commands
    a target N1, so N1 is an input, not a symptom --- which is why its deviation is zero
    by construction in our generator and why GPA treats it as the operating-point anchor.
  \item[N2 {[}rpm{]} --- core / high-pressure spool speed.] The compressor--turbine core.
    At fixed N1 and ambient, a shifting N2 re-balances the work split between spools ---
    a sensitive, low-noise indicator of where in the core something changed.
  \item[WF {[}kg/s{]} --- fuel flow.] The engine's operating cost, measured directly. At
    fixed thrust, deteriorating components demand more fuel; a 1\,\% fleet-wide WF shift
    is a fuel-bill line item, which is why airlines track it to a tenth of a percent.
  \item[EGT {[}K{]} --- exhaust gas temperature.] Measured by a ring of thermocouples
    behind the low-pressure turbine. It is the engine's \emph{thermal headroom gauge}:
    certification fixes a redline, a new engine at rated thrust on a hot day sits some
    tens of kelvin below it (the \emph{EGT margin}), and every bit of deterioration
    spends that margin. When margin reaches zero the engine must come off wing ---
    EGT margin \emph{is} the industry's remaining-life currency, which is why RUL in
    this project is defined through it.
  \item[P25, T25 {[}bar, K{]} --- booster exit / HPC inlet.] The interface between the
    low and high spools; movements here separate ``front of the engine'' from ``core''
    hypotheses.
  \item[PS3, T3 {[}bar, K{]} --- HPC exit.] The highest-pressure station, feeding the
    combustor. PS3 is the classic discriminator between compressor efficiency loss and
    flow-capacity loss --- the pair the cockpit set famously confuses
    (\cref{sec:icm}).
\end{description}

The \emph{cockpit subset} (N1, N2, WF, EGT) is what every operator has without optional
instrumentation packages; the extended set adds the four station probes. The
observability gulf between those two sets is quantified in \cref{sec:icm} and becomes an
acquisition-decision tool in \cref{ch:tomography}.

\section{What synthetic data is --- and when it is science}
\label{sec:synthetic}

\emph{Synthetic data} is data produced by a model rather than measured from the target
system~\cite{Jordon2022,Nikolenko2021}. Four families are in common use: physics-based
simulation (a mechanistic model generates observations --- this project, and NASA's
C-MAPSS/N-CMAPSS~\cite{Saxena2008,AriasChao2021}); generative statistical models (GANs,
VAEs, diffusion models fitted to real data); augmentation (perturbing real samples); and
hybrid pipelines. The motivations are always some subset of: \textbf{ground truth}
(impossible to obtain otherwise --- our case), \textbf{rarity} (failures too scarce to
learn from), \textbf{privacy/access} (airline data is commercially sensitive), and
\textbf{control} (factorial experiments no fleet would ever fly).

Synthetic data earns scientific standing only by confronting its central liability:
\emph{the model that generates the data embeds assumptions, and any method evaluated on
it might exploit those assumptions rather than solve the real problem}. This project's
defenses are layered and were built \emph{before} the headline experiments:

\begin{enumerate}
  \item \textbf{Anchor the generator to certified measurements.} The twin behind the
    generator is calibrated against type-certificate and emissions-databank measurements
    (\cref{sec:calibration}) --- its fuel-flow \emph{predictions} land within
    3\,\% of test-bed values it was never fitted to.
  \item \textbf{Separate self-consistency from fidelity.} Estimator scoring against
    ground truth is exact by construction; whether the generated physics matches the
    \emph{full nonlinear} model is an audited empirical question
    (\cref{sec:audit-nonlin}), not an assumption.
  \item \textbf{Make the dataset hard on purpose.} The difficulty gate (a deliberately
    weak classifier must \emph{fail}) exists because a synthetic dataset's most common
    disease is being accidentally easy; ours failed that gate twice and was physically
    redesigned both times (\cref{sec:audit-difficulty}).
  \item \textbf{Check the ranking externally.} The central prognosis conclusion was
    re-tested on a benchmark this project did not generate (\cref{sec:sim-to-real}).
  \item \textbf{Say what does not transfer.} Absolute performance numbers (an RMSE of
    858 cycles) are properties of this simulated world and its noise settings;
    what transfers is \emph{rankings, mechanisms and geometry} --- which method wins
    where and why, and the observability structure that any CFM56-class cockpit set
    shares. Every operational claim in this document is of the second kind.
\end{enumerate}

\section{The provenance chain, end to end}
\label{sec:provenance}

Every number in SynCFM56 descends from public, certified measurements through an
auditable chain (\cref{fig:provenance}):

\begin{figure}[htbp]
  \centering
  \begin{tikzpicture}[
      font=\footnotesize,
      box/.style={draw, rounded corners=2pt, align=center, inner sep=5pt,
                  minimum height=9mm},
      lbl/.style={font=\scriptsize\itshape, midway, above},
      >=Stealth, node distance=5mm and 9mm]
    \node[box, fill=gray!12] (src) {Certified public data\\ EASA TCDS E.004 $\cdot$ ICAO EEDB\\ \scriptsize (measured fuel flows, ratings, redlines)};
    \node[box, right=of src] (twin) {pyCycle twin\\ \scriptsize design point + off-design,\\ \scriptsize thrust-matched calibration};
    \node[box, right=of twin] (art) {F1 artifacts\\ \scriptsize baseline decks $z_0(u)$\\ \scriptsize ICM grid $\mathbf{H}(u)$};
    \node[box, below=of art] (gen) {Fleet generator\\ \scriptsize + fault catalog (YAML)\\ \scriptsize + seed 20260703};
    \node[box, left=of gen] (data) {\textbf{SynCFM56 v1.1}\\ \scriptsize snapshots $\cdot$ events $\cdot$ index\\ \scriptsize 1.55 M rows, truth attached};
    \node[box, left=of data, fill=gray!12] (use) {Consumers\\ \scriptsize audits $\cdot$ traditional $\cdot$ AI\\ \scriptsize F5/F7 verdicts};
    \draw[->] (src) -- node[lbl]{calibrate} (twin);
    \draw[->] (twin) -- node[lbl]{probe} (art);
    \draw[->] (art) -- (gen);
    \draw[->] (gen) -- node[lbl]{emit} (data);
    \draw[->] (data) -- (use);
  \end{tikzpicture}
  \caption{Data provenance: from certified test-bed measurements to the rows the
  algorithms consume. Every arrow is a versioned script; every box is a versioned
  artifact.}
  \label{fig:provenance}
\end{figure}

\begin{enumerate}
  \item \textbf{Certified measurements.} A \emph{type-certificate data sheet} is the
    regulator's public record of what an engine type demonstrated during certification:
    rated thrusts, temperature and speed limits. The \emph{ICAO Engine Emissions
    Databank} adds something quantitatively precious: sea-level test-bed \textbf{measured
    fuel flows} at the four thrust settings of the landing--takeoff certification cycle
    (100/85/30/7\,\% of rated thrust), for the exact -7B26 variant. These are real,
    traceable measurements of real hardware --- the project's only contact with physical
    reality, and deliberately its foundation.
  \item \textbf{The calibrated twin} (\cref{ch:model}) reproduces those measurements
    thrust-matched, then serves as the physics oracle.
  \item \textbf{The compressed physics}: baseline decks (healthy behavior over the
    envelope) and the ICM grid (first-order response to every health deviation) ---
    small, versioned artifacts that make fleet-scale generation tractable.
  \item \textbf{The generator} (\cref{ch:fleet}) draws each engine's degradation
    personality from the fault catalog --- the versioned YAML contract that is the
    \emph{single source of truth} for every physical and statistical assumption:
    mechanism rates, wash schedules, fault magnitudes, sensor noise, drift and dropout.
    Changing the world means editing that file and bumping the dataset version, never
    editing code.
  \item \textbf{The emitted dataset}: three tables. \code{snapshots.parquet} --- one row
    per flight cycle with the two reports; \code{events.parquet} --- the maintenance
    and fault log (washes with recovery fractions, foreign-object damage, acute episode
    onsets/parameters/magnitudes); \code{fleet_index.json} --- per-engine life, split
    assignment, severity multipliers, drift channel.
\end{enumerate}

\section{Anatomy of a SynCFM56 row}
\label{sec:row-anatomy}

Each snapshot row carries four kinds of columns, and the distinction is load-bearing:

\begin{description}
  \item[Measured channels (\code{to_*}, \code{cr_*}).] What the airline would see: truth
    passed through the sensor model --- Gaussian noise per the catalog, quantization to
    recording resolution, drift ramps on affected engines, dropout. \emph{Both EHM
    families see only these.}
  \item[True channels (\code{*_true}).] The same quantities before the sensor model ---
    kept for audits and for measuring how much each method loses to noise, never shown
    to any pipeline under evaluation.
  \item[Operating conditions (\code{to_dTs_C}, \code{cr_N1_cmd}).] The condition each
    report was taken at --- ambient temperature offset at takeoff, commanded N1 in
    cruise. Legitimate inputs (an airline knows them), and the raw material of
    \cref{ch:tomography}.
  \item[Ground truth (\code{x_*}, \code{egtm_C}, \code{rul}, labels).] The ten true
    health-parameter deviations, the true EGT margin, the true remaining useful life,
    the acute-fault isolation label, chronic condition labels and drift flags ---
    the oracle that real fleets lack, used \emph{only} for scoring and for training
    labels where the task legitimately supplies them.
\end{description}

This four-way split is the whole point of the synthetic method: the first three column
groups reproduce the epistemic situation of a real airline (noisy measurements plus known
conditions), while the fourth provides the referee that reality cannot.
